\par\SourceLeaf{435}\Needspace{7\baselineskip}\bigskip
\begin{center}
\pdfbookmark[2]{Gauss to Olbers, Göttingen, 2 December 1817}{pallas-brief-18171202}
\textsc{Gauss} to \textsc{Olbers}, Göttingen, 2 December 1817.
\end{center}
\noindent\hspace*{\parindent}. . . . .\quad I believe I have already written to you a few times about the table for making it easier to calculate the amount of Pallas's perturbations on each occasion. Since losing Mr. \textsc{Westphal}'s assistance with it, I have felt inclined to work only slowly on the part that remained, but the table has been complete for some time now. I considered it necessary to use this table to compare all 12 oppositions observed so far once again with normal elements, to derive the resulting 23 equations of condition (in one opposition the latitude is of no use, since only \textsc{David} observed declinations, and his observations were very poor, No. V. of 1808), and to determine the 7 corrections from them (including the one to the mass of {\AstroSymbols ♃}). I have now provisionally completed this laborious work (subject to a little refinement). The agreement is indeed astonishing, all the more so because I have not included the perturbations due to Saturn that Mr. \textsc{Nicolai} calculated and that would make the agreement noticeably \textls[180]{worse}. The correction to Jupiter's mass \((+\frac{1}{46})\) is now increasingly confirmed. The sum of the squares of the 12 differences in longitude, which could not be brought below \(4064''\) without this correction, is brought down to \(241''\) with it, and according to my theory of probability the determination of Jupiter's mass from the perturbations of Pallas is now about three times as accurate as that from the observations of Saturn, from which, as is well known, \textsc{Laplace} derived a result that conflicts with mine. Including \textsc{Nicolai}'s perturbations due to Saturn would, incidentally, leave the correction to Jupiter's mass unchanged. Here are the remaining differences:
\begin{center}
\begingroup\small\setlength{\tabcolsep}{5pt}
\begin{tabular}{@{}lr r r r r@{}}
&&\multicolumn{1}{c}{Heliocentr.}&\multicolumn{1}{c}{Geocentr.}&\multicolumn{2}{c}{Difference according to \textsc{Carlini}'s}\\
&&\multicolumn{1}{c}{longitude}&\multicolumn{1}{c}{latitude}&\multicolumn{2}{c}{determination of the oppositions}\\[3pt]
I.&1803&\(+0''{,}8\)&\(-1''{,}2\)&\(-5''{,}2\)&\(-30''{,}2\)\\
II.&1804&\(+4{,}0\)&\(-11{,}8\)&\(+4{,}5\)&\(-9{,}8\)\\
III.&1805&\(-2{,}1\)&\(-3{,}8\)&\(-3{,}7\)&\(+3{,}9\)\\
IV.&1807&\(-4{,}6\)&\(-2{,}3\)&\(+1{,}1\)&\(-4{,}5\)\\
V.&1808&\(-5{,}0\)&\(-16{,}5\mathpunct{::}\)&&
\end{tabular}
\endgroup
\end{center}
\SourceLeaf{436}
\begin{center}
\begingroup\small\setlength{\tabcolsep}{5pt}
\begin{tabular}{@{}lr r r r r@{}}
&&\multicolumn{1}{c}{Heliocentr.}&\multicolumn{1}{c}{Geocentr.}&\multicolumn{2}{c}{Difference according to \textsc{Carlini}'s}\\
&&\multicolumn{1}{c}{longitude}&\multicolumn{1}{c}{latitude}&\multicolumn{2}{c}{determination of the oppositions}\\[3pt]
VI.&1809&\(+0''{,}7\)&\(-5''{,}1\)&&\\
VII.&1811&\(+2{,}4\)&\(-6{,}6\)&\(+2{,}9\)&\(+6{,}0\)\\
VIII.&1812&\(+10{,}7\)&\(+4{,}8\)&&\\
IX.&1813&\(-0{,}7\mathpunct{\raisebox{1ex}{.}}\)&\(-0{,}8\)&&\\
X.&1814&\(+1{,}8\)&\(-1{,}1\)&&\\
XI.&1816&\(-2{,}7\)&\(-8{,}5\)&&\\
XII.&1817&\(-6{,}5\)&\(+5{,}0\)&&
\end{tabular}
\endgroup
\end{center}
\par These differences are all the less significant because the determinations of several oppositions by other astronomers, e.g. \textsc{Carlini}, often differ from mine by more than the differences above, As appetite comes with eating, I intend to have oppositions II and VIII calculated once again from the observations using my osculating elements. I also intend to have the observations of 1802 analyzed. The longitudes now agree even better than the latitudes, which is no surprise, since the former are heliocentric and the latter geocentric.
\par
